
\noindent Nesta seção, são apresentados os principais recursos computacionais utilizados no
desenvolvimento deste trabalho, incluindo ferramentas para escrita do documento,
linguagens de programação empregadas na implementação dos modelos e bibliotecas utilizadas
na resolução dos problemas de otimização convexa. A escolha dessas ferramentas foi guiada
pela necessidade de garantir rigor científico, reprodutibilidade dos experimentos e
eficiência computacional na implementação dos algoritmos propostos.

\section{Ambiente de desenvolvimento e escrita}

A redação deste trabalho foi realizada utilizando o sistema de preparação de documentos
\LaTeX, amplamente adotado na comunidade científica devido à sua capacidade de lidar de
forma robusta com expressões matemáticas, referências bibliográficas e formatação de
documentos técnicos. O uso do \LaTeX\ foi particularmente importante neste trabalho,
dado o elevado número de equações e demonstrações envolvidas na formulação dos modelos
baseados em máquinas de vetores de suporte.

Para a edição do código-fonte e do documento, foi utilizado o editor Visual Studio Code
(VSCode), que oferece um ambiente leve e altamente configurável. O VSCode permite integração
com diversas extensões, possibilitando suporte eficiente tanto para a escrita em \LaTeX\
quanto para o desenvolvimento em linguagens de programação como Python. Entre as
funcionalidades utilizadas, destacam-se a navegação facilitada entre arquivos,
realce de sintaxe, execução de scripts e integração com sistemas de controle de versão.

\section{Linguagem de programação e implementação dos modelos}

A implementação dos modelos propostos neste trabalho foi realizada na linguagem de
programação Python, devido à sua ampla adoção na área de ciência de dados e aprendizado
de máquina, bem como à disponibilidade de bibliotecas eficientes para computação numérica.

O Python foi utilizado tanto para a geração de dados simulados quanto para a implementação
dos algoritmos de Support Vector Regression (SVR) e sua extensão para dados censurados
(SVCR e WSVCR). A manipulação de dados foi realizada com o auxílio da biblioteca NumPy,
que oferece estruturas eficientes para operações matriciais e vetoriais, fundamentais
para a implementação de métodos baseados em kernel.

Além disso, o Python permite uma integração natural com bibliotecas de otimização e
facilita a prototipagem de modelos matemáticos complexos, o que foi essencial para o
desenvolvimento e validação das abordagens propostas neste trabalho.

\subsection{Otimização convexa}

A resolução dos problemas de otimização associados aos modelos SVCR e WSVCR foi realizada
por meio da biblioteca CVXOPT, uma biblioteca especializada em programação convexa
desenvolvida para a linguagem Python.

A escolha do CVXOPT deve-se ao fato de que os problemas dual dos modelos baseados em
máquinas de vetores de suporte podem ser formulados como problemas de programação
quadrática convexa. Nesse contexto, o CVXOPT fornece rotinas eficientes e numericamente
estáveis para a obtenção da solução ótima global, garantindo a correta resolução do
problema sem a necessidade de implementação manual de algoritmos iterativos, como o
Sequential Minimal Optimization (SMO).

O uso de uma biblioteca consolidada como o CVXOPT também contribui para a reprodutibilidade
dos resultados, uma vez que os algoritmos de otimização empregados são amplamente testados
e documentados na literatura.
